\chapter{The Theory: Purity, Monotonicity, Convergence}
\label{ch:theory}

The architecture chapters established \emph{how} \caem{} works. This chapter asks
\emph{why it should work at all}, and answers the question \Cref{ch:problem} left open:
if the verifier is imperfect, why does the memory not slowly fill with the verifier's
mistakes? The answer is a small set of results about the \emph{stored pool}. We teach
each one from the ground up, state precisely what it does and does not claim, and pair
it with what the realised run actually showed.

\begin{intuition}[The one question the theory must answer]
Every gate in \caem{} is fallible. The verifier is better than chance but not perfect,
so it will sometimes admit a wrong answer. The natural worry is a feedback spiral: the
verifier's errors enter memory, the loop trains on them, the model degrades, and the
errors compound. The theory's job is to show that this does not happen, that under a
mild condition the stored pool ends up \emph{cleaner} than the verifier itself, and that
the system contracts toward a stable, high-purity state rather than spiralling. The
whole chapter is the unpacking of that promise.
\end{intuition}

\section{What the theorems are about, and what they are not}
\label{sec:theory-scope}

Before any result, one scope statement, because it governs how every claim should be
read.

\begin{pitfall}[These are pool-purity guarantees, not model-accuracy guarantees]
The theorems bound the \emph{purity of the stored memory pool}: the fraction of stored
episodes whose answer is correct. They do \emph{not} bound the exact-match accuracy of
the fine-tuned model. The two are different quantities flowing through different
channels: the pool purity is a property of what the gate admits, while the model's
accuracy is a property of what gradient descent does with the training pool. A cycle in
which the model's accuracy on some benchmark dips is therefore \emph{not} a violation of
these theorems; it is simply outside their scope. Keeping this distinction sharp is what
lets the theory make clean, provable claims about the memory while the messier question
of the trained model's accuracy is handled empirically in \Cref{ch:trajectory}. When this
chapter says ``purity,'' it always means the stored pool, never the model.
\end{pitfall}

\begin{precise}[Pool purity, defined]
Let $\mathcal{M}_t$ be the set of episodes in memory at the end of cycle $t$. Its
\emph{purity} is
\[
  \pi_t \;:=\; \Pr\big(\text{answer correct} \,\big|\, \text{episode} \in \mathcal{M}_t\big),
\]
the chance that a randomly drawn stored episode is right. Purity one means a spotless
pool; purity equal to the base model's accuracy means the gate added nothing. The
theory is the study of how $\pi_t$ behaves across cycles.
\end{precise}

\section{The data-purity result: why an imperfect verifier is enough}
\label{sec:theory-purity}

This is the heart of the chapter and the direct answer to \Cref{ch:problem}'s objection.
It rests on two quantities. The \emph{base accuracy} $p$ is how often the generator is
right before any filtering. The \emph{balanced accuracy} $\alpha$ is how good the gate is
at telling right from wrong, averaged over the two cases, with $\alpha = 1/2$ meaning a
coin flip and $\alpha = 1$ a perfect judge.

\begin{precise}[The purity identity]
If the gate admits an answer, what is the chance it is actually correct? By Bayes' rule,
writing the gate's true-positive and true-negative rates both as $\alpha$ for clarity,
\[
  \pi \;=\; \frac{p\,\alpha}{p\,\alpha \;+\; (1-p)(1-\alpha)}.
\]
Read the denominator as ``ways to be admitted'': a correct answer admitted (rate
$p\,\alpha$) plus a wrong answer wrongly admitted (rate $(1-p)(1-\alpha)$). The purity is
the first term's share. The crucial algebraic fact is that whenever $\alpha > 1/2$, this
$\pi$ is strictly greater than $p$: \textbf{the pool is purer than the unfiltered
generator}, and for a strong enough generator it is purer than the verifier's own
accuracy. A gate only modestly better than a coin flip still manufactures a clean pool.
\end{precise}

\begin{intuition}[Why a so-so judge produces a clean pool]
The trick is that the gate is applied to a \emph{population}, and base rates do the heavy
lifting. Suppose the generator is right $60\%$ of the time and the gate is right $70\%$ of
the time. Among the answers the gate admits, the correct ones were both common to begin
with and likely to be passed, while the wrong ones were less common and likely to be
caught. The arithmetic above turns those two advantages into an admitted-pool purity of
about $78\%$, higher than either the generator's $60\%$ or the gate's $70\%$. This is the
base-rate argument \Cref{ch:problem} promised, and it is why \caem{} does not need a
perfect verifier, only one that beats a coin flip on the claims it scores.
\end{intuition}

\begin{figure}[t]
\centering
\begin{tikzpicture}
\begin{axis}[
  width=10.6cm, height=6.6cm,
  xlabel={base accuracy $p$ (generator correct rate)},
  ylabel={pool purity $\pi$},
  xmin=0, xmax=1, ymin=0, ymax=1,
  xtick={0,0.25,0.5,0.75,1}, ytick={0,0.25,0.5,0.75,1},
  axis lines=left, font=\small\sffamily, tick label style={font=\scriptsize},
  legend style={at={(0.97,0.05)}, anchor=south east, draw=cbInk!25, font=\scriptsize, fill=white},
  grid=both, grid style={cbInk!8}, domain=0:1, samples=120,
]
  \addplot[cbInk!45, dashed, line width=1pt] {x};
  \addlegendentry{$\pi = p$ (no gate)}
  \addplot[cbIntuition, line width=1.4pt] {0.6*x/(0.6*x + 0.4*(1-x))};
  \addlegendentry{$\alpha = 0.6$}
  \addplot[cbScenario, line width=1.4pt] {0.7*x/(0.7*x + 0.3*(1-x))};
  \addlegendentry{$\alpha = 0.7$}
  \addplot[cbExample!70!black, line width=1.4pt] {0.8*x/(0.8*x + 0.2*(1-x))};
  \addlegendentry{$\alpha = 0.8$}
  \addplot[cbScenario!60!cbInk, line width=1.4pt, densely dotted] {0.67*x/(0.67*x + 0.33*(1-x))};
  \addlegendentry{$\alpha \approx 0.67$ (deployed)}
  \addplot[cbExample, only marks, mark=*, mark size=2.6pt] coordinates {(0.54,0.71)};
  \addlegendentry{deployed point}
\end{axis}
\end{tikzpicture}
\caption[The purity identity]{Pool purity $\pi$ as a function of base accuracy $p$, for
three gate qualities $\alpha$. Every curve lies \emph{above} the diagonal $\pi = p$
whenever $\alpha > 1/2$: the gate makes the pool purer than the unfiltered generator. A
gate that is only modestly better than chance ($\alpha = 0.6$) still lifts a
$60\%$-accurate generator to a pool purity near $0.69$, and a $\alpha = 0.7$ gate lifts it
to about $0.78$, above the gate's own accuracy. The dotted curve at $\alpha \approx 0.67$
and the marked operating point are the deployed cold-start gate: at the realised base rate
$p \approx 0.54$ the identity predicts a pool purity near $0.71$, which sits just under the
realised calibration-fold purity band of $0.72$ to $0.75$, so the deployed run lands in the
band the identity anticipates. (Curve family illustrative; the marked point and the
$\alpha \approx 0.67$ curve are from the deployed cold-start fold.)}
\label{fig:purity}
\end{figure}

\subsection*{The cross-cycle floor}

The identity is the expected purity for one admission. The first theorem turns it into a
guarantee that holds \emph{every} cycle, accounting for the fact that the pool is the
union of newly admitted episodes and old episodes that survived re-verification.

\begin{precise}[Pool purity invariant]
Let $\pi_{\text{store}}$ be the gate's calibration-fold precision at the storage
threshold, and $\pi_{\text{retro}}$ the precision of episodes that survive the
retroactive re-verification prune. Then for every cycle,
\[
  \pi_{t+1} \;\ge\; \min\!\big(\pi_{\text{store}},\, \pi_{\text{retro}}\big).
\]
The reasoning is a one-line convexity argument: next cycle's pool is the disjoint union
of survivors (each at least $\pi_{\text{retro}}$ pure) and new admissions (each at least
$\pi_{\text{store}}$ pure), and the purity of a union sits between the two, hence above
the smaller. The bound is one-sided and can be re-checked from each cycle's own
calibration fold, which makes it a falsifiable contract rather than an aspiration.
\end{precise}

\begin{bythenumbers}[The floor, and what the run showed]
The registered floor is $\pi_{\text{store}} \approx 0.64$ at the storage threshold
$\taustore = 0.60$, with the retroactive prune at $\tauretro = 0.50$ supplying the second
floor. Across the realised six cycles the pooled calibration-fold purity stayed in the
band $0.72$ to $0.75$, clearing the floor by eight to eleven points at every cycle. The
guarantee held, with margin to spare. A reading below the floor would have falsified
either the exchangeability assumption or the prune contract; none was observed.
\end{bythenumbers}

\begin{pitfall}[The floor is a pooled guarantee, and one benchmark dips below it]
The floor is stated on the \emph{pooled} training-panel axis, not per benchmark, because
the assumption it rests on (that the calibration fold and the candidate pool are
interchangeable) holds on the pooled axis but not necessarily on each benchmark's slice.
On the realised run the FEVER per-benchmark reading dips just below $0.64$ at cycles one
through four (in the band $0.61$ to $0.63$) and recovers at cycle five. This is reported
transparently because it is real, but it is not a falsification: it is a per-benchmark
diagnostic outside the pooled scope the theorem actually claims. Stating the scope
precisely is what lets the dip be informative rather than alarming.
\end{pitfall}

\section{Monotonicity: a better generator stores more}
\label{sec:theory-monotone}

The second theorem connects the loop's effect on the generator to the gate's behaviour.
As the self-improvement loop sharpens the generator, the calibration-fold separation
between correct and incorrect answers, measured by Cohen's $d$ (a standard effect size:
the gap between the two groups' means in units of their spread), widens. The theorem asks
what that does to the storage rate $\sigma$, the fraction of answers admitted.

\begin{precise}[Sign of the storage-rate response]
Modelling the storage score as approximately Gaussian within each correctness class, the
direction in which the storage rate moves as the separation $d$ grows is
\[
  \operatorname{sign}\!\Big(\frac{\partial \sigma}{\partial d}\Big)
  \;=\; \operatorname{sign}\big(p_+\phi_+ - p_-\phi_-\big),
\]
where $\phi_\pm$ are the densities of the correct and incorrect classes at the threshold.
In words: when the operating point sits in the high-precision regime, where the correct
class has more mass at the threshold than the incorrect class, a better generator (larger
$d$) admits \emph{more} answers, $\sigma_{t+1} \ge \sigma_t$. Improving the model and
holding the threshold fixed lets more of its now-better answers clear the bar.
\end{precise}

\begin{pitfall}[The realised receipt is honest about being partial]
This is the theorem whose empirical support is weakest, and the thesis says so. Of the
three observable cycle-to-cycle transitions, one aligns strongly with the prediction, one
is too small to read, and one moves the opposite way. The reverse-sign transition is not
treated as a refutation but is attributed to two effects the simple Gaussian model leaves
out: the composite is re-fitted each cycle (so the storage score is not a fixed function),
and the within-class spread contracts across cycles (so the shared-variance assumption is
approximate). A clean test would freeze the composite for a cycle and let only the
generator move, which is registered as future work. Reporting a load-bearing theorem with
a partial receipt, and naming exactly why, is the honest posture.
\end{pitfall}

\section{Convergence: the gap contracts to a stable band}
\label{sec:theory-convergence}

The third theorem describes the long-run trajectory of purity. Writing $\pi_\infty$ for
the asymptotic purity and $G_t = |\pi_t - \pi_\infty|$ for the remaining gap, the gate-plus-loop
update behaves as a contraction.

\begin{precise}[Contraction envelope, and why it is reframed]
Under a stationarity assumption on the per-cycle storage rate and separation, the gap
shrinks geometrically in expectation,
\[
  \mathbb{E}[G_t] \;\le\; (1-r)^t\, G_0, \qquad r \in (0,1).
\]
Each cycle admits fresh high-purity episodes and re-scores old ones under the improved
generator, and the expected purity moves a fixed fraction $r$ of the way toward its limit.
But six cycles is far too short to estimate a rate $r$ cleanly, so the claim the thesis
actually registers is the weaker, testable one: \emph{bounded-band stationarity}, that
$\pi_t$ stays inside a narrow band rather than diverging. On the realised run the
calibration-fold purity traced a band of width below $0.03$, consistent with contraction
and inconsistent with a spiral; the sharper geometric-rate test is deferred to a longer
horizon.
\end{precise}

\begin{intuition}[The idealised limit, and why reality sits below it]
A companion identity shows that, in an idealised world where the verifier's quality stays
fixed above one half and nothing else interferes, the purity recurrence climbs strictly
toward one: keep filtering with a better-than-chance gate and the pool tends to perfection.
\caem{} does not reach that ideal, and the theory says why. Three deliberate regularisers
pin the realised ceiling strictly below one: the frozen backbone and bounded adapter cap
how far the model can move, the retention guard turns any regressive cycle into a no-op, and
the finite capacity of the generator and the retrieval corpus impose a hard floor on error.
The gap between the realised plateau and the theoretical one is the price of these safety
mechanisms, not a failure of the mathematics.
\end{intuition}

\section{The Tier-1 floor: memory hallucination is bounded by pool impurity}
\label{sec:theory-tier1}

The first corollary cashes the purity guarantee out into a user-facing rate. The
memory-direct tier serves a stored answer verbatim, so its error rate is exactly the
pool's impurity. Hence
\[
  \limsup_{t\to\infty}\, \mathrm{CHM}_1(t) \;\le\; \max\!\big(1-\pi_{\text{store}},\, 1-\pi_{\text{retro}}\big),
\]
where $\mathrm{CHM}_1$ is the memory tier's hallucination metric and $\limsup$ is the
long-run worst case (it tolerates bounded wobble rather than demanding a smooth descent).

\begin{pitfall}[A capability bound, not a trajectory]
This corollary must be read as a statement of \emph{capability}, not a measured curve. By
design the evaluation set is held disjoint from the training stream, so the memory tier
fires on almost nothing during evaluation, seven times across nine thousand held-out
queries, every one returning the correct stored answer. The asymptotic hallucination rate
of a tier that rarely fires is not observable on this horizon. So the corollary says ``the
memory tier cannot hallucinate faster than the pool is impure,'' a sound bound, and the
realised receipt is the targeted routing probe rather than a trajectory. The conservative
ceiling of about $0.36$ from the cold-start fold comfortably overshoots the realised pool
impurity of roughly $0.25$, which the retroactive prune tightens further.
\end{pitfall}

\section{Self-correction: bad episodes are pruned in finite time}
\label{sec:theory-selfcorrect}

The most reassuring corollary addresses the feedback-spiral worry head on. Each stored
episode is re-scored at every cycle boundary by the same verifier against the improved
generator. Treat surviving a re-verification pass as a noisy yes/no test of the episode's
correctness.

\begin{precise}[Why a wrong episode cannot survive forever]
For a truly correct episode, each pass is more likely to keep it than drop it; for a wrong
one, more likely to drop it than keep it, provided the verifier's balanced accuracy stays
above one half. The evidence from repeated passes \emph{multiplies}: the likelihood ratio
in favour of correctness, given survival of $k$ passes, grows without bound, so the chance
a survivor is actually correct tends to one. Concretely, a wrong episode's per-cycle prune
probability is at least $\delta = 2\alpha - 1 > 0$, so its expected survival time is at most
$1/\delta$ cycles, a finite number. Hallucinations do not accumulate; they are flushed.
\end{precise}

\begin{figure}[t]
\centering
\begin{tikzpicture}
\begin{axis}[
  width=10.4cm, height=5.6cm,
  xlabel={re-verification passes survived $k$}, ylabel={chance a wrong episode survives},
  xmin=0, xmax=8, ymin=0, ymax=1, xtick={0,2,4,6,8}, ytick={0,0.25,0.5,0.75,1},
  axis lines=left, font=\small\sffamily, tick label style={font=\scriptsize},
  legend style={at={(0.97,0.97)}, anchor=north east, draw=cbInk!25, font=\scriptsize, fill=white},
  grid=both, grid style={cbInk!8}, domain=0:8, samples=9,
]
  \addplot[cbPitfall, line width=1.4pt, mark=*, mark size=1.2pt] {0.8^x};
  \addlegendentry{$\alpha=0.6$ \,($\delta=0.2$)}
  \addplot[cbNumbers!75!black, line width=1.4pt, mark=square*, mark size=1.2pt] {0.6^x};
  \addlegendentry{$\alpha=0.7$ \,($\delta=0.4$)}
  \addplot[cbExample!70!black, line width=1.4pt, mark=triangle*, mark size=1.4pt] {0.4^x};
  \addlegendentry{$\alpha=0.8$ \,($\delta=0.6$)}
\end{axis}
\end{tikzpicture}
\caption[Self-correction]{(Illustrative.) A wrong episode's chance of surviving $k$
re-verification passes decays geometrically, at a rate set by the verifier's margin over
chance $\delta = 2\alpha-1$. Even a modest gate ($\alpha=0.6$) drives the survival chance
below ten percent within ten passes; a stronger gate flushes errors in a handful. This is
why repeated re-verification prevents a feedback spiral rather than causing one. The
deployed run carries no per-episode $k$-pass survival record; the directional receipt that
the chapter does report, the per-cycle prune rate declining from about thirteen to four
percent, is a distinct aggregate quantity rather than a measurement of these curves.}
\label{fig:selfcorrect}
\end{figure}

\begin{bythenumbers}[The directional receipt]
The clean prediction (every wrong episode pruned eventually) is asymptotic and not testable
in six cycles, so the registered receipt is directional: the per-cycle prune rate should
decline as the pool's quality rises. It did, monotonically, from thirteen percent at cycle
one to nine, then seven, then four percent at cycle five, as the verifier found fewer
low-quality entries to remove. The trend is the supported reading; an exact rate is deferred
to a longer run.
\end{bythenumbers}

\section{Stochastic equilibrium: learning under a guard that sometimes says no}
\label{sec:theory-equilibrium}

The last corollary folds the retention guard of \Cref{ch:cycle} into the convergence picture.
Because a cycle commits only when the worst retention probe holds above the floor, the commit
events form a coin-flip sequence with some success rate $\pi_{\text{commit}}$, and the
contraction runs only on the cycles that commit. Three consequences follow.

\begin{description}[leftmargin=0pt, labelsep=0.5em, style=nextline]
  \item[The commit rate is a real fraction.] On any trajectory where the guard ever fires,
  $\pi_{\text{commit}}$ lies strictly between zero and one. The realised sequence of commit
  decisions was keep, keep, keep, abort, keep, a commit fraction of four in five, with
  worst-probe ratios $1.006$, $0.949$, $0.976$, $0.886$, $0.958$ against the floor $0.93$.
  \item[Calendar progress slows, but does not stop.] The contraction envelope on the calendar
  timeline becomes $(1 - \pi_{\text{commit}}\cdot r)^t$: aborted cycles simply do not advance
  the model, so convergence is paced by the commit rate.
  \item[Memory grows regardless.] The fresh-data path runs on every cycle, committed or not,
  so the memory and the per-tier coverage grow monotonically even across an abort. This is the
  asymmetric rollback of \Cref{ch:cycle} stated as an invariant: weights advance only on commits,
  but memory advances always.
\end{description}

\section{The results, assembled}
\label{sec:theory-map}

The results lock together, as \Cref{fig:theory-map} shows. The purity identity and its
cross-cycle floor establish that the pool stays clean; monotonicity describes how the gate
responds as the generator improves; convergence bounds the trajectory to a stable band;
the Tier-1 floor turns pool purity into a user-facing rate; self-correction guarantees errors
are flushed; and stochastic equilibrium accounts for the guard. Together they answer
\Cref{ch:problem}'s objection in full: an imperfect verifier is enough, errors do not
compound, and the system contracts toward a stable, high-purity state.

\begin{figure}[t]
\centering
\begin{tikzpicture}[
  font=\small\sffamily,
  th/.style={rectangle, rounded corners=4pt, draw=cbFormula!70!black, fill=cbFormula!8,
    line width=1pt, align=center, inner sep=4pt, text width=3.0cm, font=\footnotesize, minimum height=0.9cm},
  cor/.style={rectangle, rounded corners=4pt, draw=cbExample!70!black, fill=cbExample!8,
    line width=1pt, align=center, inner sep=4pt, text width=3.0cm, font=\footnotesize, minimum height=0.9cm},
  flow/.style={-{Stealth[length=2.4mm]}, line width=1pt, draw=cbInk!70},
]
  \node[th] (bayes) at (0,1.4) {\textbf{purity identity}\\{\footnotesize $\alpha>1/2 \Rightarrow \pi>p$}};
  \node[th] (floor) at (4.2,1.4) {\textbf{purity floor}\\{\footnotesize $\pi_{t+1}\ge\min(\cdot,\cdot)$}};
  \node[th] (mono) at (0,-0.4) {\textbf{monotonicity}\\{\footnotesize better gen.\ stores more}};
  \node[th] (conv) at (4.2,-0.4) {\textbf{convergence}\\{\footnotesize bounded band}};
  \node[cor] (tier1) at (8.6,1.4) {\textbf{Tier-1 floor}\\{\footnotesize $\mathrm{CHM}_1\le 1-\pi$}};
  \node[cor] (self) at (8.6,0.0) {\textbf{self-correction}\\{\footnotesize errors pruned in $\le 1/\delta$}};
  \node[cor] (eq) at (8.6,-1.4) {\textbf{stochastic equilibrium}\\{\footnotesize commit-paced, memory grows}};
  \draw[flow] (bayes) -- (floor);
  \draw[flow] (floor) -- (tier1);
  \draw[flow] (mono) -- (conv);
  \draw[flow] (conv) -- (eq);
  \draw[flow] (floor) -- (self);
  \draw[flow] (conv) -- (tier1);
\end{tikzpicture}
\caption[How the results lock together]{The dependency map. The purity identity (an
imperfect gate yields a clean pool) underpins the cross-cycle floor; the floor and the
convergence bound feed the Tier-1 hallucination ceiling and the finite-time self-correction;
monotonicity and convergence, qualified by the retention guard, give the stochastic
equilibrium. Purple nodes are the three load-bearing theorems; green are the corollaries. (Schematic.)}
\label{fig:theory-map}
\end{figure}

\begin{defence}[If an examiner asks: ``do these theorems prove CAEM reduces hallucination?'']
They prove a precise, bounded thing and not a sweeping one. They establish that the stored
\emph{memory pool} stays clean under a verifier only modestly better than chance, that it
cannot spiral, and that the memory tier's error is bounded by the pool's impurity. They do
\emph{not} prove that the fine-tuned model's accuracy rises on every benchmark; that is a
property of the training gradient, outside the pool-level scope, and is settled empirically
in \Cref{ch:trajectory}. The honest framing is that the theory secures the data foundation,
a clean, self-correcting memory, and the empirical chapters measure what the model built on
that foundation actually achieves. Several results carry partial or asymptotic receipts at
the six-cycle horizon, and each is reported with its scope rather than overclaimed.
\end{defence}

\begin{bythenumbers}[The theory in constants]
\begin{itemize}[leftmargin=1.3em, itemsep=1pt]
  \item Precondition: verifier balanced accuracy $\alpha > 1/2$ on the claims it scores.
  \item Purity floor: $\min(\pi_{\text{store}}, \pi_{\text{retro}})$, registered at $\approx 0.64$; realised band $0.72$ to $0.75$.
  \item Tier-1 hallucination ceiling: $\max(1-\pi_{\text{store}}, 1-\pi_{\text{retro}})$; realised pool impurity $\approx 0.25$.
  \item Self-correction: wrong episode expected survival $\le 1/\delta$ cycles, $\delta = 2\alpha - 1$; prune rate fell $13\% \to 4\%$.
  \item Commit fraction: $\pi_{\text{commit}} = 4/5$ realised; convergence paced by $(1-\pi_{\text{commit}} r)^t$.
  \item Scope: all results bound the \emph{stored pool}, not the model's exact-match accuracy.
\end{itemize}
\end{bythenumbers}

The theory secures the foundation. \Cref{ch:benchmarks} turns to measurement: the five
benchmarks, the eight external baselines, and the metrics, the exact-match and
composite-hallucination axes, that the next chapters use to report what the realised system
achieved.
